\section{Problem Formulation}
\label{sec:problem}

\paragraph{Dispatch and completion} The agent observes an environment state $s_t$ through an observation $\Omega(s_t)$ (for the browser: URL, title, element list and screenshot). A planner maps a goal $g$, a plan $P=(g_1,\dots,g_n)$ and the observation to an action $a_t$. Executing it returns a \emph{dispatch flag} $\alpha_t\in\{0,1\}$, which says only that the call returned without error. The agent may report the task complete only when a completion predicate holds on the state observed \emph{after} the action:
\begin{equation}
\Phi(g,s_{t+1},P,\sigma) = \textstyle\bigwedge_{c\in C(g,P)} c(s_{t+1}) \;\wedge\; \pi(P) \;\wedge\; \neg\beta(s_{t+1},\sigma),
\label{eq:phi-general}
\end{equation}
where $C(g,P)$ is the set of postconditions that the goal and plan imply, $\pi(P)$ states that no plan step is pending, and $\beta$ is a blocking condition (a bot-check page, an unapproved high-risk action) with $\sigma$ the approval state. Dispatch does not imply completion, $\alpha_t=1\not\Rightarrow\Phi=1$. Every evaluation of $\Phi$ is stored with the observed values it used. The same form applies to desktop, file and process state, whose postconditions are typed and exact (Appendix~\ref{app:computer}); only the browser instance is implemented.

\paragraph{The implemented browser predicate} $C(g,P)$ is derived from the wording of the request. $\Phi$ is the conjunction of: no browser-error URL, error or ``not found'' title, or redirect to a sign-in page the request did not name ($E$); not still on a search engine's home page when a search, an extraction or a multi-step plan was requested ($H$); a results page carrying the query quoted in $g$ ($Q$); the literal text that $g$ asks to be entered present in an input, exactly if $g$ says ``exactly'' ($T$); for extraction requests, an answer or structured extraction that mentions a content word of $g$ ($X$); and, for multi-step plans, no pending step unless a structured extraction exists ($R$). Each conjunct is instantiated only when its trigger words occur, so for a request that triggers none, $\Phi$ reduces to $\neg E$. A pre-check implements $\beta$ for bot-check pages and halts the task as \textsc{Blocked}. A plan step counts as done when an action other than \texttt{complete} dispatches during it. $\Phi$ has no tunable threshold: its operating point is fixed by the rules, so we characterize the trade-off between false completions and false rejections by removing conjuncts (Section~\ref{sec:res-ablation}) rather than by a curve.

\paragraph{Metrics} Let $\mathcal{T}$ be a set of episodes or end states, $\mathcal{T}_c\subseteq\mathcal{T}$ those reported complete, and $\Phi^{\star}$ an oracle written independently of $\Phi$ (Section~\ref{sec:method}). With $N_{\mathrm{FC}}=\lvert\{\tau\in\mathcal{T}_c:\Phi^{\star}(\tau)=0\}\rvert$ and $N_{\mathrm{FR}}=\lvert\{\tau\in\mathcal{T}\setminus\mathcal{T}_c:\Phi^{\star}(\tau)=1\}\rvert$,
\begin{equation}
\mathrm{FCR} = \frac{N_{\mathrm{FC}}}{\lvert\mathcal{T}_c\rvert}, \quad
\mathrm{FRR} = \frac{N_{\mathrm{FR}}}{\lvert\{\tau\in\mathcal{T}:\Phi^{\star}(\tau)=1\}\rvert}.
\label{eq:fcr}
\end{equation}
Because a dispatch-only rule reports every state complete, we also report the share of unsatisfied states that the gate withholds, $1-N_{\mathrm{FC}}/\lvert\{\tau:\Phi^{\star}(\tau)=0\}\rvert$. FCR alone rewards an agent that never finishes, so the rates are always reported together, as counts with Wilson 95\% intervals~\cite{wilson1927probable}.
